\documentclass{article}
\usepackage{fullpage,amsmath,amsthm,graphicx,enumitem,amssymb}
\usepackage[hidelinks]{hyperref}
\usepackage{todonotes}
\usepackage{float}
\usepackage{tikz}


\begin{document}


  \section{The Linear Gaussian Information Filter (IF)}
    Modifies the classical KF to address numerical issues with inverses and posdef requirements.
    Define \( I^{-}_{k+1} = [P^{-}_{k+1}]^{-1}\) Predicted information matrix. \newline
       \( I^{+}_{k+1} = [P^{+}_{k+1}]^{-1}\) -$>$ Updated information matrix.

       Helps primarily when p (measurments) $>$ n. Won't need to invert innovation matrix.
       Can also initialize with no information.

       Caveats: Need to reconstruct the actual state estimate. 
       \[
          \hat{x}^+_{k+1} = [I^+_{k+1}]^{-1} \hat{i}^+_{k+1}
       \]

       Requires $I$ to be invertible, this means the system must be observable. If you need 3
       measurments for the system to be observable, you won't be able to reconstruct a state estimate
       unitl you observe three measurments.

       \subsection{When to use IF vs KF}

      IF more computationally efficient if p$>$n and if $R_k$ is diagonal. (More sensors than states).
      



    
\end{document}
